\begin{center}
\(\ast\qquad\ast\)
\end{center}

We thus arrive at what may be called the golden age of geometry, which falls roughly between the publication dates of Monge's Descriptive Geometry (1795) (X) and F. Klein's "Erlangen Program" (1872) (XXV b). The essential advances we owe to this sudden revival of geometry are the following:

\begin{enumerate}
\def\labelenumi{\Alph{enumi})}
\item
  The notion of an element at infinity (point, line, or plane), introduced by Desargues in the seventeenth century (VI), but which hardly manifests itself in the eighteenth century only as an abuse of language, is rehabilitated and systematically used by Poncelet (XII), who thus makes projective space the general framework for all geometric phenomena.
\item
  At the same time, with Monge, and above all Poncelet, the transition to complex projective geometry takes place. The notion of imaginary point, sporadically used during the eighteenth century, is exploited here (concurrently with that of the point at infinity) to yield statements independent of the "cases of figure" of real affine geometry. If, at first, the justifications offered in support of these innovations remain highly awkward (especially on the part of the adherents of the "synthetic" school of geometry, where the use of coordinates came to be regarded as a stain), one cannot fail to recognize therein, under the name of "principle of contingent relations" in Monge, or of "principle of continuity" in Poncelet, the first germ of the idea of "specialization" in modern algebraic geometry (*).
\end{enumerate}

One of the first results arising from these conceptions is the observation that, in complex projective space, all nondegenerate conics (resp. quadrics) are of the same nature; this leads Poncelet to the discovery of the "isotropic" elements: "Circles placed arbitrarily in a plane," he says, "are therefore not entirely independent of one another, as one might at first believe; they have ideally two imaginary points in common at infinity" ((XII), p.~48). Further on, he likewise introduces the "umbilic," the imaginary conic at infinity common to all spheres ((XII), p.~370); and if he does not speak particularly of the isotropic generators of the sphere, he at least explicitly emphasizes the existence of rectilinear generators, real or imaginary, for all quadrics (ibid., p.~371) (*); notions of which his successors (notably Plücker and Chasles), even more than he himself, make great use, in particular in the study of the "focal" properties of conics and quadrics.

\begin{enumerate}
\def\labelenumi{\Alph{enumi})}
\setcounter{enumi}{2}
\item
  The notions of point transformation and composition of transformations are likewise formulated in a general way and introduced systematically as means of proof. Apart from displacements and projections, only a few particular transformations had been known up to that time: certain planar projective transformations, of the type \(x' = a/x\) , \(y' = y/x\) used by La Hire and Newton, the “affinity” \(x' = ax\) , \(y' = by\) from Clairaut and Euler, and finally certain particular quadratic transformations, in Newton again, Maclaurin, and Braikenridge. Monge, in his Descriptive Geometry, shows all the use that can be made of plane projections in three-dimensional geometry. In Poncelet, one of the systematic methods of proof, employed to excess, consists in reducing by projection the properties of conics to those of the circle (a method already applied on occasion by Desargues and Pascal); and in order to pass similarly from a quadric to a sphere, he invents the first example of a projective transformation in space, the "homology" ((XII), p.~357); finally, it is also he who introduces the first examples of birational transformations of a curve into itself. In 1827, Möbius ((XIII a), p.~217) (and independently Chasles in 1830 (XV b)) define the most general projective linear transformations; at the same period there appear inversion (cf.~§ 10, exerc. 13) and other types of quadratic transformations, whose study will inaugurate the theory of birational transformations, which will develop in the second half of the nineteenth century.
\item
  The notion of duality appears in full light and is consciously linked to the theory of bilinear forms. The theory of poles and polars with respect to conics, which since Apollonius had made progress only in the work of Desargues and La Hire, is extended to quadrics by Monge, who, together with his students, perceives the possibility of transforming known theorems into new results by this means (*). But it is again to Poncelet that the credit belongs for having erected these remarks into a general method in his theory of transformations "by reciprocal polars," and for having made of it a particularly effective tool of discovery. A little later, notably with Gergonne, Plücker, Möbius, and Chasles, the general notion of duality detaches itself from the connection with quadratic forms, still too narrow in Poncelet. In particular, Möbius, in examining the various possibilities of duality in 3-dimensional space (defined by a bilinear form), discovers in 1833 the duality with respect to an alternating bilinear form (XIII b) (**), especially studied in the nineteenth century, in the form of the theory of “linear complexes” (cf.~§ 10, exerc. 16) and developed in connection with the “geometry of lines” and the “Plücker coordinates” introduced by Cayley, Grassmann, and Plücker around 1860.
\item
  From the very beginnings of projective geometry, the intensive study of the properties of classical geometry in their relations with projective space had quickly led to dividing them into "projective properties" and "metric properties"; and it is doubtless no exaggeration to see in this separation one of the clearest manifestations, at that time, of what was to become the modern notion of structure. But Poncelet, who first introduced this distinction and this terminology, was already aware of what connects these two types of properties; and, addressing in his Treatise the problems concerning angles, whose properties "do not seem to form part of those we have called projective\ldots, they nevertheless follow in so simple a manner," he said, "from the principles that form the basis {[}of this work{]}\ldots, which I do not believe that any other geometric theory can lead to in a manner at once more direct and simpler. One will not be at all surprised by this, if one considers that the projective properties of figures are necessarily the most general of those that can belong to them; so that they must include, as simple corollaries, all the other particular properties or relations of extension" ((XII), p.~248). To tell the truth, after this declaration, one is somewhat surprised to see him approach questions of angles in a very roundabout way, by attaching them to the focal properties of conics, instead of bringing the cyclic points directly into play; and in fact, it was only 30 years later that Laguerre (still a student at the École Polytechnique) gave the expression of an angle between lines by means of the cross-ratio of these lines and the isotropic lines with the same origin (cf.~§ 10, exerc. 5) ((XXI), vol. II, p.~13). Finally, with Cayley (XVIII d) there is clearly expressed the fundamental idea that the "metric" properties of a plane figure are none other than the "projective" properties of the figure augmented by the cyclic points --- a decisive milestone toward the "Erlangen program".
\item
  Hyperbolic non-Euclidean geometry, which came into being around 1830, at first remains somewhat apart from the movement whose main lines we are tracing. Arising from concerns of an essentially logical order touching on the foundations of classical geometry, this new geometry is presented by its inventors (*) in the same axiomatic and "synthetic" form as Euclid's geometry, and without any connection to projective geometry (whose introduction following the classical model even seemed excluded a priori, since the notion of a unique parallel disappears in this geometry); it is doubtless for this reason that for a long time it attracted little interest from the French, German, and English schools of projective geometry. Thus, when Cayley, in the fundamental memoir cited above (XVIII d), had the idea of replacing the cyclic points (considered as a conic "degenerate tangentially") by an arbitrary conic (which he called the "absolute"), he in no way thought of connecting this idea to the geometry of Lobachevsky-Bolyai, although he indicates how his conception leads to new expressions for the "distance" between two points, and although he mentions its links with spherical geometry. The situation changes around 1870, when the non-Euclidean geometries, following the diffusion of the works of Lobachevsky, and the publication of the works of Gauss and of Riemann's inaugural lecture, came to the forefront of mathematical actuality. Following the path traced by Riemann, Beltrami, without knowing Cayley's work, rediscovered in 1868 the expressions for distance given by the latter, but in a quite different context, by considering the interior of a circle as an image of a surface of constant curvature, in which geodesics are represented by straight lines (XXIV); it is Klein who, two years later, made (independently of Beltrami) the synthesis of these various points of view, which he completed by the discovery of elliptic non-Euclidean space (XXV a) (**).
\item
  In the second half of the period we are considering here, a period of critical reflection sets in, during which the partisans of "synthetic" geometry, not content with having banished coordinates from their proofs, claim to do without real numbers even in the axioms of geometry. The principal representative of this school is von Staudt, who essentially succeeded in performing this tour de force (XXIII), much admired in his time and even well into the twentieth century; and if today one no longer attributes the same importance to ideas of this order, whose possibilities for fruitful application have proved rather meager, one must nevertheless recognize that the efforts of von Staudt and his disciples contributed to clarifying ideas on the role of real or complex "scalars" in classical geometry, and thereby to introducing the modern conception of geometries over an arbitrary base field.
\end{enumerate}

Around 1860, "synthetic" geometry was at its zenith, but the end of its reign was fast approaching. Having remained heavy and inelegant throughout the eighteenth century, analytic geometry, in the hands of Lamé, Bobillier, Cauchy, Plücker, and Möbius, finally acquired the elegance and concision that would enable it to compete on equal terms with its rival. Above all, from about 1850 onward, the ideas of group and invariant, at last formulated precisely, gradually took center stage, and it became apparent that the theorems of classical geometry are nothing other than the expression of identical relations among invariants or covariants of the similarity group (*), just as those of projective geometry express the identities (or "syzygies") among covariants of the projective group. This is the thesis masterfully expounded by F. Klein in the celebrated "Erlangen program" (XXV b), where he advocates abandoning the sterile controversies between the "synthetic" and "analytic" tendencies; if, he says, the accusation leveled against the latter of granting a privileged role to an arbitrary system of axes "was only too often justified with regard to the defective manner in which the coordinate method was formerly used, it collapses when it comes to a rational application of this method.\ldots{} The domain of spatial intuition is not forbidden to the analytic method.\ldots{}", and he emphasizes that “one must not underestimate the advantage that a well-suited formalism brings to subsequent research, in that it anticipates thought, so to speak” ((XXV b), pp. 488–490).

This leads to a rational and "structural" classification of the theorems of "geometry" according to the group from which they derive: the linear group for projective geometry, the orthogonal group for metric questions, and the symplectic group for the geometry of the "linear complex." But under this pitiless clarity, classical geometry—with the exceptions of algebraic geometry and differential geometry (**), now constituted as autonomous sciences—suddenly withers and loses all its brilliance. Already the generalization of methods based on the use of transformations had rendered the formation of new theorems somewhat mechanical: "Today," said Chasles in 1837 in his \emph{Aperçu historique}, "anyone may come forward, take any known truth, and subject it to the various general principles of transformation; he will extract from it other truths, different or more general; and these will be susceptible to similar operations; so that one may multiply, almost to infinity, the number of new truths deduced from the first.\ldots{} Anyone who wishes may therefore, in the current state of science, generalize and create in Geometry; genius is no longer indispensable for adding a stone to the edifice" ((XV b), p.~268-269). But the situation becomes much clearer with the progress of invariant theory, which finally succeeds (at least for the “classical” groups) in formulating general methods that make it possible in principle to write out all algebraic covariants and all their “syzygies” in a purely automatic manner; a victory which, at the same time, marks the death, as a field of research, of classical invariant theory itself, and of “elementary” geometry, which has practically become a mere dictionary for it. To be sure, nothing allows one to foresee a priori, among the infinity of “theorems” that can thus be unrolled at will, which ones will have, in an appropriate geometric language, a simplicity and elegance comparable to the classical results, and there remains a restricted domain where numerous amateurs continue to exercise themselves happily (geometry of the triangle, of the tetrahedron, of algebraic curves and surfaces of low degree, etc.). But for the professional mathematician, the mine is exhausted, since there are no longer any problems of structure there, capable of having repercussions on other parts of mathematics; and this chapter of group theory and invariant theory may be considered closed until further notice. (*)

\[
\ast \ast
\]

Thus, after the Erlangen program, Euclidean and non-Euclidean geometries, from a purely algebraic point of view, became simple languages, more or less convenient, for expressing the results of the theory of bilinear forms, whose progress goes hand in hand with that of the theory of invariants (*). Everything concerning the notion of the rank of a bilinear form and the relations between these forms and linear transformations is definitively clarified by the work of Frobenius (XXVII a). It is also to Frobenius that we owe the canonical expression of an alternating form on a free Z-module (§ 5, no. 1, th. 1) (XXVII b); however, skew-symmetric determinants had already appeared in Pfaff, at the beginning of the century, in connection with the reduction of differential forms to a normal form; Jacobi, who in 1827 takes up this problem again (XIX a), knows that a skew-symmetric determinant of odd order is zero, and it is he who forms the expression of the Pfaffian and shows that it is a factor of the skew-symmetric determinant of even order; but he had not perceived that the latter is the square of the Pfaffian, and this point was established only by Cayley in 1849 (XVIII b). The notion of the symmetric bilinear form associated with a quadratic form is the most elementary case of the process of "polarization," one of the fundamental tools of the theory of invariants. Under the name of "scalar product," this notion would enjoy immense fortune, first with the popularizers of "vector calculus," then, from the twentieth century, thanks to the unforeseen generalization brought to it by the theory of Hilbert space (see the historical note in Book V). It is also this latter theory that will bring to light the notion of the adjoint of an operator (which previously had scarcely manifested itself except in the theory of linear differential equations, and, in tensor calculus, through the dance of co- and contravariant indices under the baton of the metric tensor); it is this theory, finally, that will give full prominence to the notion of Hermitian form, introduced first by Hermite in 1853 in connection with arithmetic investigations ((XXII), p.~237), but which remained somewhat on the margins of the main mathematical currents until around 1925 and the applications of complex Hilbert spaces to quantum theories.

The study of the orthogonal group and the group of similarities --- clearly conceived and treated as such since the middle of the nineteenth century, and having become the heart of the theory of quadratic forms — as well as of the other "classical" groups (the linear group, the symplectic group, and the unitary group) — takes on, on the other hand, an increasingly important role. We can only mention here the essential part played by these groups in the theory of Lie groups and differential geometry on the one hand, and in the arithmetic theory of quadratic forms (see, for example, (XXXIII) and (XXXV)) on the other (**); to this circumstance, as well as to the extension of the concept of duality to the most diverse questions, is due the fact that there is scarcely any modern mathematical theory in which bilinear forms do not intervene in one way or another. We must in any case note that it was the study of the rotation group (in three dimensions) that led Hamilton to the discovery of quaternions (XVII); this discovery was generalized by W. Clifford, who in 1876 introduced the algebras bearing his name and proved that they are tensor products of quaternion algebras, or of quaternion algebras and a quadratic extension (XXVIII). Rediscovered four years later by Lipschitz (XXIX), who used them to give a parametric representation of orthogonal transformations in n variables (generalizing those that Cayley had obtained for n = 3 (XVIII a) and n = 4 (XVIII c) via the theory of quaternions (cf. § 9, exercises 15 and 16)), these algebras, and the notion of "spinor" derived from them (see (XXXII) and (XXXIV)), were also to enjoy great popularity in the modern era by virtue of their use in quantum theories.

Finally, we have only a word left to say about the evolution of ideas that led to the almost total abandonment of any restriction on the ring of scalars in the theory of sesquilinear forms — a tendency common to all modern algebra, but one that perhaps manifested itself here earlier than elsewhere. We have already noted the fruitful introduction of geometry over the field of complex numbers (which, moreover, throughout the nineteenth century, was not without a perpetual and sometimes perilous confusion between this geometry and real geometry); here the clarity comes above all from the axiomatic studies of the end of the nineteenth century on the foundations of geometry (XXX). In the course of these investigations, Hilbert and his followers, in particular, while examining the relations among the various axioms, were led to construct suitable counterexamples, where the “base field” (commutative or not) possessed more or less pathological properties, and they thus accustomed mathematicians to “geometries” of an entirely new type. From the analytic point of view, Galois had already considered linear transformations where coefficients and variables took their values in a finite prime field ((XVI), p.~27); in developing these ideas, Jordan (XXVI) is naturally led to envisage the classical groups over these fields, groups whose intervention manifests itself in varied domains of mathematics. Dickson, around 1900, extended Jordan's investigations to all finite fields, and more recently it has been realized that a large part of the Jordan–Dickson theory extends to the case of an absolutely arbitrary “base field”; this is due essentially to the general properties of isotropic vectors and to Witt's theorem, which, trivial in the classical cases, were established for an arbitrary base field only in 1936 (XXXI) (*).

But in thus pushing the study of sesquilinear forms toward an ever greater “abstraction,” it proved extremely suggestive to retain as is the terminology which, in the case of 2- and 3-dimensional spaces, came from classical geometry, and to extend it to the n-dimensional case and even to infinite-dimensional spaces. Outmoded as an autonomous and living science, classical geometry was thereby transfigured into a universal language of contemporary mathematics, of incomparable flexibility and convenience.

\section*{Bibliography}
\phantomsection
\addcontentsline{toc}{section}{Bibliography}

\begin{enumerate}
\def\labelenumi{(\Roman{enumi})}
\item
  O. NEUGEBAUER, Lectures on the History of Ancient Mathematical Sciences, Vol. I: Pre-Greek Mathematics, Berlin (Springer), 1934.
\item
  J. TROPFKE, History of Elementary Mathematics, vols. IV–VI, Berlin–Leipzig (de Gruyter), 1923–24.
\item
  Euclid's Elements, 5 vols., ed. J. L. Heiberg, Leipzig (Teubner), 1883–88.
\end{enumerate}

(III bis) T. L. HEATH, The thirteen books of Euclid's Elements\ldots, 3 vols., Cambridge, 1908.

\begin{enumerate}
\def\labelenumi{(\Roman{enumi})}
\setcounter{enumi}{3}
\item
  T. L. HEATH, Apollonius of Perga, Treatise on conic sections, Cambridge (Univ. Press), 1896.
\item
  Ptolemaei Cl. Opera, ed. J. L. Heiberg, 2 vols., Leipzig (Teubner), 1898–1903.
\item
  G. DESARGUES, WORKS\ldots vol. I, Paris (Leiber), 1864: Brouillon project d'une atteinte aux événements des rencontres d'un cône avec un plan, pp. 103-230.
\item
  P. FERMAT, Œuvres, vol. I, Paris (Gauthier-Villars), 1891: a) Ad locos planos et solidos Isagoge, pp. 91-110 (French translation, ibid., vol. III, pp. 84-101); b) Isagoge ad locos ad superficiem, pp. 111-117 (French translation, ibid., vol. III, pp. 102-108).
\item
  L. EULER: a) Introductio in Analysin Infinitorum (Opera Omnia (1), vol. IX, Zürich-Leipzig-Berlin (O. Füssli and B. G. Teubner), 1945); b) Theoria motus corporum solidorum seu rigidorum (Opera Omnia (2), vol. III, Zürich-Leipzig-Berlin (O. Füssli and B. G. Teubner), 1948); c) Problema algebraicum ob affections prorsus singulares memorabile (Opera Omnia, (1), vol. VI, Leipzig-Berlin (Teubner), 1921, pp. 287-315); d) Formulae generales pro translatione quacunque corporum rigidorum, Novi Comm. Acad. Sc. imp. Petrop., vol. XX (1776), pp. 189-207; e) De centro similitudinis. (Opera Omnia (1), vol. XXVI, Zürich (O. Füssli), 1956, pp. 276-285).
\item
  J. L. LAGRANGE, Œuvres, Paris (Gauthier-Villars), 1867-1892: a) Researches on the method of maxima and minima, vol. I, pp. 3–20; b) Researches in arithmetic, vol. III, pp. 695–795.
\item
  G. MONGE, Descriptive Geometry, Paris, 1798.
\item
  C. F. GAUSS, Works: a) Disquisitiones arithmeticae, vol. I, Göttingen, 1870; b) Mutations of Space, vol. VIII, Göttingen, 1900, pp. 357–362.
\item
  J.-V. PONCELET, Treatise on the Projective Properties of Figures, vol. I, 2nd ed., Paris (Gauthier-Villars), 1865.
\item
  A. F. Möbius, Collected Works, Leipzig (Hirzel), 1885–87: a) The Barycentric Calculus, vol. I, pp. 1–388; b) On a Special Kind of Dual Relations between Figures in Space, vol. I, pp. 489–515 (= Crelle's Journal, vol. X, 1833); c) On a New Treatment of Analytic Spherics, vol. II, pp. 1–54.
\item
  A.-L. CAUCHY: a) Lessons on the applications of infinitesimal calculus to geometry (Complete Works, (2), vol. V, Paris (Gauthier-Villars), 1903); b) On the equation by means of which one determines the secular inequalities of the planets (Complete Works, (2), vol. IX, Paris (Gauthier-Villars), 1891, pp. 174–195).
\item
  M. CHASLES: a) Note on the general properties of the system of two bodies, Bull. de Férussac, vol. XIV (1830), pp. 321–326; b) Historical survey of the origin and development of methods in geometry, Brussels, 1837.
\item
  E. GALOIS, Mathematical Works, Paris (Gauthier-Villars), 1897.
\item
  W. R. HAMILTON, Lectures on Quaternions, Dublin, 1853.
\item
  A. CAYLEY, Collected Mathematical Papers, Cambridge, 1889-1898: a) On certain results relating to quaternions, vol. I, p.~123-126 (= Phil. Mag., 1845); b) Sur les déterminants gauches, vol. I, p.~410-413 (= J. de Crelle, vol. XXXVIII (1848)); c) Recherches ultérieures sur les déterminants gauches, vol. II, p.~202-215 (= J. de Crelle, vol. L (1855)); d) A sixth memoir on quantics, vol. II, p.~561-592 (= Phil. Trans., 1859).
\item
  C. G. J. JACOBI, Gesammelte Werke, Berlin (G. Reimer), 1881-1891: a) On the Pfaffian method\ldots vol. IV, pp. 17–29; b) On a fundamental algebraic theorem and its applications, vol. III, pp. 593–598.
\item
  J. J. SYLVESTER, Collected Mathematical Papers, vol. I, Cambridge, 1904: A demonstration of the theorem that every homogeneous quadratic polynomial is reducible by real orthogonal substitution to the form of a sum of positive and negative squares, pp. 378–381 (= Phil. Mag., 1852).
\item
  E. LAGUERRE, Works, vol. II, Paris (Gauthier-Villars), 1905.
\item
  C. HERMITE, Works, vol. I, Paris (Gauthier-Villars), 1905: On the theory of quadratic forms, pp. 200–263 (= J. de Crelle, vol. XLVII (1854)).
\item
  K. G. V. von STÄUDT, Contributions to the Geometry of Position, Nuremberg, 1856.
\item
  E. BELTRAMI: a) Essay on the interpretation of non-Euclidean geometry, Giorn. di Mat., vol. VI (1868), pp. 284–312; b) Fundamental theory of spaces of constant curvature, Ann. di Mat. (2), vol. II (1868–69), pp. 232–255.
\item
  F. KLEIN, Collected Mathematical Works, vol. I, Berlin (Springer), 1921: a) On the so-called non-Euclidean geometry, p.~254-305 (= Math. Ann., vol. IV (1871)); b) Comparative considerations on recent geometrical research, p.~460-497 (= Math. Ann., vol. XLIII (1893)).
\item
  C. JORDAN, Treatise on substitutions and algebraic equations, Paris (Gauthier-Villars), 1870.
\item
  G. FROBENIUS; a) On linear substitutions and bilinear forms, J. de Crelle, vol. LXXXIV (1878), p.~1-63; b) Theory of linear forms with integer coefficients, J. de Crelle, vol. LXXXVI (1879), p.~146-208.
\item
  W. K. CLIFFORD, Mathematical Papers, London (Macmillan), 1882: a) On the classification of geometric algebras, p.~397-401; b) Applications of Grassmann's extensive algebras, p.~266-276 (= Amer. Journ. of Math., vol. I (1878)).
\item
  R. LIPSCHITZ, Investigations on the sums of squares, Bonn, 1886.
\item
  D. HILBERT, Foundations of Geometry, Leipzig (Teubner), 1899.
\item
  E. WITT, Theory of quadratic forms in arbitrary fields, J. de Crelle, vol. CLXXVI (1937), p.~31-44.
\item
  E. CARTAN, Lectures on the theory of spinors, Actual. Sci. et Industr., nos. 643 and 701, Paris (Hermann), 1938.
\item
  C. L. SIEGEL, Symplectic Geometry, Amer. Journ. of Math., vol. LXV (1943), p.~1-86.
\item
  C. CHEVALLEY, The algebraic theory of spinors, New York (Columbia Univ. Press), 1954.
\item
  M. EICHLER, Quadratic forms and orthogonal groups, Berlin-Göttingen-Heidelberg (Springer), 1952.
\end{enumerate}

\chapter*{Index of Notations}
\phantomsection
\addcontentsline{toc}{chapter}{Index of Notations}

\(b^{J}, b^{J'} : 1, 2.\)

The reference numbers indicate successively the section and the number (or, exceptionally, the exercise).

\(F^{J}\) (an antiautomorphism of the ring B of scalars of the right module F): 1, 2.

\(N^{0}, M^{0}\) (N, M submodules): 1, 3.

\(u^{*}\) (u homomorphism): 1, 8.

\(M(x), \mathbf{x} (x \text{ element of a free module}) : 1, 10.\)

\(M(u)\) (u homomorphism from a free module to a free module): 1, 10.

\(^{t}M, \dot{M}^{J} (M \text { matrix } ): 1, 10.\)

\(\mathrm{D}_{\Phi}(x_{1},\ldots ,x_{n}),\mathrm{D}_{\Phi}(\mathrm{S}):2.\)

\(\overline{\alpha}\) ( \(\alpha\) scalar): 3.

\(\operatorname{Pf}(R):5,2.\)

\(\mathbf{Sp}(\Phi), \mathbf{Sp}(2m, \mathrm{A}), \mathbf{Sp}_{2m}(\mathrm{A}) : 5, 3.\)

\(\mathbf{U}(\Phi), \mathbf{S}\mathbf{U}(\Phi) (\Phi \text{ Hermitian form}) : 6, 2.\)

O(Q), SO(Q) (Q quadratic form): 6, 2.

\(\mathbf{U}(n,\mathrm{A}),\mathbf{S}\mathbf{U}(n,\mathrm{A}),\mathbf{O}(n,\mathrm{A}),\mathbf{SO}(n,\mathrm{A}):6,2.\)

\(Q^{\top}Q^{\prime}, Q \sim Q^{\prime}(Q, Q^{\prime}\text{ quadratic forms}): 8, 1.\)

θ(Q) (Q quadratic form): 8, 1.

\(\mathbf{T} + \mathbf{T}', a. \mathbf{T} (\mathbf{T}, \mathbf{T}'\) types of quadratic forms): 8, 2.

\(\delta (\mathbb{Q})\) (Q quadratic form): 8, 2.

\(\mathbf{Q} \otimes \mathbf{Q}'\) (Q, \(\mathbf{Q}'\) quadratic forms): 8, 3.

\(\mathrm{TT}^{\prime}\) (T, \(\mathrm{T}^{\prime}\) types of quadratic forms): 8, 3.

\(C(Q), I(Q), \rho_{Q}, \rho, C^{+}(Q), C^{-}(Q), C^{+}, C^{-} (Q \text{ quadratic form}): 9, 1.\)

\(\alpha, \beta, \mathrm{G}(f): 9, 1.\)

\(e_x, i_f, i_x^{\mathrm{F}}, \lambda_{\mathrm{F}}: 9, 2.\)

\(\lambda_{\mathbf{F}}:9,3.\)

A(Φ) (Φ symmetric bilinear form on a vector space of dimension 2): 10, 1.

\(\overline{u} (u\) direct similarity) \(i,d:10,1\) .

\((\widehat{\mathbf{D}_1,\mathbf{D}_2})\) \((\mathbf{D}_1,\mathbf{D}_2\) lines or half-lines): 10, 3.

\[
\mathfrak {A}, \mathfrak {A} _ {0}, h, h ^ {\prime}: 1 0, 3.
\]

\(|x|, \widehat{(x,y)}(x,y\text{ vectors}):10,3.\)

cos θ, sin θ, tan θ, cot θ: 10, 3.

\(\left.\right)D_{1}, D_{2}\left\{\left(D_{1}, D_{2}\right)\left(D_{1}, D_{2} \text{ half-lines in an oriented plane}\right): 10, 4.\right.\)

\chapter*{Terminological Index}
\phantomsection
\addcontentsline{toc}{chapter}{Terminological Index}

The reference numbers indicate successively the paragraph and the number (or, exceptionally, the exercise).

Adjoint (right, left) of a homomorphism: 1, 8.

Adjoint of a semilinear map: 7, 3.

Clifford algebra of a quadratic form: 9, 1.

Alternating (matrix): 3, 1.

Angle of two half-lines: 10, 3, and exercise 6.

Angle between two rays in an affine plane: 10, 3.

Angle between two lines: 10, 3 and exercise 2.

Angle of two lines in an affine plane: 10, 3.

Angle between two directed lines: 10, exercise 2.

Angle between two vectors in a plane: 10, 3.

Angle between two vectors in a space of dimension \(\geqslant 2:10,3\) .

Right angle: 10, 3 and exercise 2.

Angle of a rotation: 10, 3.

Angle of a direct similarity: 10, exercise 6.

Straight angle: 10, 3 and exercise 6.

Ring of types of quadratic forms: 8, 3.

Witt ring: 8, 3.

Antiautomorphism of a ring: 1, 2.

Principal antiautomorphism of a Clifford algebra: 9, 1.

Antihermitian (form): 3, 1.

Antihermitian (matrix): 3, 1.

Antisymmetric (matrix): 3, 1.

Bilinear map: 1, 1.

Degenerate bilinear map (on the right, on the left): 1, 1.

Non-degenerate bilinear map: 1, 1.

Nondegenerate bilinear map associated with a bilinear map: 1, 3.

Bilinear map obtained by extension of scalars from a bilinear map: 1, 4.

Canonical map of a module into its Clifford algebra: 9, 1.

Canonical mapping from the set of Hermitian endomorphisms to the set of Hermitian forms: 7, 3.

Canonical mapping of the set \(\mathfrak{A}\) on \(\mathrm{S}^{+}/\mathrm{H}:10,3\) .

Canonical mapping of the set \(\mathfrak{A}\) on \(O^{+}:10,3\) .

Right (left) linear map associated with a bilinear map: 1, 1.

Right (left) semilinear map associated with a sesquilinear form: 1, 6.

Right (left) sesquilinear map for an antiautomorphism: 1, 2. Degenerate (right, left) sesquilinear map: 1, 2. Nondegenerate sesquilinear map: 1, 2. Nondegenerate sesquilinear map associated with a sesquilinear map: 1, 3. Sesquilinear map obtained by extension of scalars from a sesquilinear map: 1, 4. Associated (linear map) with a bilinear map: 1, 1. Associated (semilinear map) with a sesquilinear form: 1, 6. Associated (bilinear form) with a quadratic form: 3, 4. Orthogonal automorphism: 6, 2. Orthogonal automorphism in n variables: 6, 2. Principal automorphism of a Clifford algebra: 9, 1. Symplectic automorphism: 5, 3. Symplectic automorphism in 2m variables: 5, 3. Unitary automorphism: 6, 2. Unitary automorphism in n variables: 6, 2. Axis of an affine linear complex: 10, exerc. 16. Orthogonal basis: 6, 1. Orthonormal basis: 6, 1. Symplectic basis: 5, 1. Bilinear (map): 1, 1. Bilinear (form): 1, 6. Bimodule: 1, 1. Canonical: see Canonical map and Canonical homomorphism. Center of a quadric: 6, exerc. 25 and 10, exerc. 12. Unit circle: 10, exerc. 2. Chasles (relation of): 10, 3. Clifford (algebra of): 9, 1. Clifford (group of): 9, 5. Linear complex: 10, exerc. 16. Condition (C): 6, 1. Condition (C'): 6, 1. Condition (T): 4, 2. Isotropic cone: 6, exerc. 23. Conformal (group): 10, exerc. 14. Affine conic: 6, exerc. 25. Projective conic: 6, exerc. 23. Conjugate (affine linear varieties): 6, exerc. 25. Conjugate (projective linear varieties): 6, exerc. 23. Cosine of an angle: 10, 3. Cotangent of an angle: 10, 3. Witt decomposition: 4, 2. Degenerate (bilinear map) on the right (left): 1, 1. Degenerate (sesquilinear map) on the right (left): 1, 2. Degenerate (affine conic): 6, exerc. 25. Degenerate (projective conic): 6, exerc. 23. Degenerate (affine quadric): 6, exerc. 25. Degenerate (projective quadric): 6, exerc. 23. Half-line: 10, 3. Closed half-line: 10, 3. Isotropic half-line: 10, 3. Open half-line: 10, 3. Displacement: 6, 6. Dimension of a quadratic form: 8, 1. Direct (similarity): 6, 5. Discriminant of a sesquilinear form with respect to a system of elements: 2. Distance: 7, 1. Right (angle): 10, 3. Pointed line: 10, exerc. 2.

Odd element of a Clifford algebra: 9, 1. Isotropic element: 4, 1. Even element of a Clifford algebra: 9, 1. Singular element: 4, 1. Elements orthogonal with respect to a sesquilinear map: 1, 3. Elements orthogonal with respect to a quadratic form: 3, 4. Hermitian endomorphism: 7, 3. Normal endomorphism: 7, 3. ε-Hermitian (form): 3, 1. ε-Hermitian (matrix): 3, 1. Equivalent (quadratic forms): 3, 4. Equivalent (sesquilinear forms): 1, 6. Definition space of a quadratic form: 8, 1. Euclidean space: 6, 6. Hermitian space: 6, 6. Oriented space: 10, 3. Euclidean (space): 6, 6. Extension of a sesquilinear form to a tensor power: 1, 9. Extension of a sesquilinear form to an exterior power: 1, 9. External (direct sum) of quadratic forms: 3, 4. External (direct sum) of quadratic modules: 3, 4. Weakly orthogonal (subspaces): 3, exerc. 11. Closed (angular sector): 10, 4. Closed (half-line): 10, 3. Cosine function: 10, 3. Cotangent function: 10, 3. Sine function: 10, 3. Tangent function: 10, 3. Trigonometric function: 10, 3. Anti-Hermitian form: 3, 1. Bilinear form: 1, 6. Bilinear form associated with a quadratic form: 3, 4. Hermitian bilinear form: 3, 1. Hermitian form: 3, 1. Negative (positive) Hermitian form: 7, 1. Inverse form of a bilinear (sesquilinear) form: 1, 7. Metric form: 6, 6. Neutral form: 4, 2 and 8, 1. Quadratic form: 3, 4. Negative (positive) quadratic form: 7, 1. Non-degenerate quadratic form: 3, 4. Quadratic form obtained by extension of scalars from a quadratic form: 3, 4. Sesquilinear form: 1, 6. Equivalent quadratic forms: 3, 4. Equivalent sesquilinear forms: 1, 6. Gram-Schmidt (orthogonalization process of): 6, 1. Clifford group: 9, 5. Reduced Clifford group: 9, 5. Special Clifford group: 9, 5. Rotation group: 9, 5. Group of types of quadratic forms: 8, 2. Witt group: 8, 2. Orthogonal group associated with Q: 6, 2. Orthogonal group in n variables: 6, 2. Special orthogonal group: 6, 2. Special unitary group: 6, 2. Symplectic group associated with Φ: 5, 3. Symplectic group in 2m variables: 5, 3. Unitary group associated with Φ: 6, 2.

\begin{Shaded}
\begin{Highlighting}[]
\NormalTok{Groupe unitaire à n variables : 6, 2.}
\NormalTok{Hermitien (endomorphisme) : 7, 3.}
\NormalTok{Hermitien (espace) : 6, 6.}
\NormalTok{Hermitienne (forme) : 3, 1.}
\NormalTok{Hermitienne (matrice) : 3, 1.}
\NormalTok{Homomorphisme canonique de l\textquotesingle{}algèbre de Clifford d\textquotesingle{}un sous{-}module de E dans l\textquotesingle{}algèbre de Clifford de E : 9, 1.}
\NormalTok{Homomorphisme métrique : 4, 3.}
\NormalTok{Hyperplan radical de deux sphères : 10, exerc. 12.}
\NormalTok{Image réciproque d\textquotesingle{}une application bilinéaire : 1, 1.}
\NormalTok{Image réciproque d\textquotesingle{}une application sesquilinéaire : 1, 2.}
\NormalTok{Impair (élément) dans une algèbre de Clifford : 9, 1.}
\NormalTok{Indice d\textquotesingle{}une forme quadratique (sesquilinéaire) : 4, 2.}
\NormalTok{Invariant de Dickson : 9, exerc. 9.}
\NormalTok{Inverse (forme) : 1, 7.}
\NormalTok{Inverse (similitude) : 6, 5.}
\NormalTok{Inversion : 10, exerc. 13.}
\NormalTok{Inversion de sphère C : 10, exerc. 13.}
\NormalTok{Involution dans un groupe linéaire : 6, 3.}
\NormalTok{Isotrope (demi{-}droite) : 10, 3.}
\NormalTok{Isotrope (élément) : 4, 1.}
\NormalTok{Isotrope (sous{-}module) : 4, 1.}
\NormalTok{Isotrope (variété linéaire) : 6, 6.}
\NormalTok{Laguerre (formule de) : 10, exerc. 5.}
\NormalTok{Loi d\textquotesingle{}inertie : 7, 2.}
\NormalTok{Longueur d\textquotesingle{}un vecteur : 10, 3.}
\NormalTok{Matrice alternée : 3, 1.}
\NormalTok{Matrice antihermitienne : 3, 1.}
\NormalTok{Matrice antisymétrique : 3, 1.}
\NormalTok{Matrice d\textquotesingle{}une application satisfaisant à (G), (D), (GD) ou (DG) : 1, 10.}
\NormalTok{Matrice d\textquotesingle{}une application bilinéaire : 1, 1.}
\NormalTok{Matrice d\textquotesingle{}une application sesquilinéaire : 1, 2.}
\NormalTok{Matrice ε{-}hermitienne : 3, 1.}
\NormalTok{Matrice hermitienne : 3, 1.}
\NormalTok{Matrice normale : 7, exerc. 17.}
\NormalTok{Matrice orthogonale : 6, 2.}
\NormalTok{Matrice symétrique : 3, 1.}
\NormalTok{Matrice symplectique : 5, 3.}
\NormalTok{Matrice unitaire : 6, 2.}
\NormalTok{Métrique (forme) : 6, 6.}
\NormalTok{Métrique (homomorphisme) : 4, 3.}
\NormalTok{Module quadratique : 3, 4.}
\NormalTok{Multiplicateur d\textquotesingle{}une similitude : 6, 5 et 6.}
\NormalTok{Négatif (n{-}vecteur) : 10, 3.}
\NormalTok{Négative (forme hermitienne) : 7, 1.}
\NormalTok{Négative (forme quadratique) : 7, 1.}
\NormalTok{Neutre (forme quadratique) : 8, 1.}
\NormalTok{Neutre (forme sesquilinéaire) : 4, 2.}
\NormalTok{Non dégénérée (application bilinéaire) : 1, 1.}
\NormalTok{Non dégénérée (application sesquilinéaire) : 1, 2.}
\NormalTok{Non dégénérée (forme quadratique) : 3, 4.}
\NormalTok{Normal (endomorphisme) : 7, 3.}
\NormalTok{Norme spinorielle : 9, 5.}
\NormalTok{Noyau d\textquotesingle{}un module quadratique : 3, 4.}
\NormalTok{n{-}vecteur négatif (positif) : 10, 3.}
\NormalTok{Orientation d\textquotesingle{}un espace : 10, 3.}
\NormalTok{Orienté (espace) : 10, 3.}
\NormalTok{Orthogonal (automorphisme) : 6, 2.}
\NormalTok{Orthogonal (groupe) : 6, 2.}
\NormalTok{Orthogonal (groupe spécial) : 6, 2.}
\end{Highlighting}
\end{Shaded}

Orthogonal (projector): 6, 3. Orthogonal (submodule) to a submodule: 1, 3. Orthogonal (vector) to a linear manifold: 6, 6. Orthogonal (basis): 6, 1. Orthogonal (matrix): 6, 2. Orthogonal (projection) onto a linear manifold: 6, 6. Orthogonal (transformation): 6, 2. Orthogonal (manifold) to a linear manifold: 6, 6. Orthogonal (parts): 1, 3 and 3, 4. Orthogonal (linear manifolds): 6, 6. Orthogonalization (process of) of Gram-Schmidt: 6, 1. Orthogonal (elements): 1, 3 and 3, 4. Orthonormal (basis): 6, 1. Open (angular sector): 10, 4. Open (half-line): 10, 3. Even (element) in a Clifford algebra: 9, 1. Orthogonal parts: 1, 3 and 3, 4. Perpendicular (linear manifolds): 6, exerc. 22. Pfaffian of an alternating matrix: 5, 2. Flat (angle): 10, 3. Point of view of a stereographic projection: 10, exerc. 14. Polar of a linear manifold with respect to a quadric: 6, exerc. 23 and 25. Pole of a hyperplane with respect to a quadric: 6, exerc. 23 and 25. Pole of an inversion: 10, exerc. 13. Positive (n-vector): 10, 3. Positive (Hermitian form): 7, 1. Positive (quadratic form): 7, 1. Principal (antiautomorphism): 9, 1. Principal (automorphism): 9, 1. Tensor product of quadratic forms: 8, 3. Tensor product of sesquilinear forms: 1, 9. Orthogonal projector: 6, 3. Orthogonal projection onto a linear manifold: 6, 6. Stereographic projection: 10, exerc. 14. Pseudo-discriminant: 9, exerc. 9. Power of a point with respect to a sphere: 10, exerc. 12. Power of an inversion: 10, exerc. 13. Pythagoras (theorem of): 6, 6. Quadratic (form): 3, 4. Quadratic (module): 3, 4. Affine quadric: 6, exerc. 25. Projective quadric: 6, exerc. 23. Rank of a bilinear (sesquilinear) form: 1, 6. Radius of a sphere: 10, exerc. 12. Reduced (Clifford group): 9, 5. Chasles relation: 10, 3. Spinor representations: 9, 4. Rotation: 9, 5. Rotation of angle φ: 10, 3 and exerc. 2. Rotations (group of): 9, 5. Closed angular sector: 10, 4. Open angular sector: 10, 4. Flat angular sector: 10, exerc. 8. Reentrant angular sector: 10, exerc. 8. Salient angular sector: 10, exerc. 8. Semi-spinor: 9, 4. Sesquilinear (map): 1, 2. Sesquilinear (form): 1, 6. Signature of a Hermitian form: 7, 2. Similarity: 6, 5 and 6.

Direct similarity: 6, 5 and 9, exercise 9.

Inverse similarity: 6, 5.

Singular (element): 4, 1.

Singular (submodule): 4, 1.

Sine of an angle: 10, 3.

Direct sum of bilinear (sesquilinear) maps: 1, 3.

External direct sum of quadratic forms (quadratic modules): 3, 4.

Isotropic submodule: 4, 1.

Orthogonal submodule to a submodule: 1, 3.

Singular submodule: 4, 1.

Totally isotropic submodule: 4, 1.

Totally singular submodule: 4, 1.

Sphere: 10, exercise 12.

Orthogonal spheres: 10, exercise 12.

Spinor: 9, 4.

Direct sequence of half-lines: 10, 4.

Symmetry with respect to a hyperplane: 6, 4.

Symmetry with respect to a vector subspace: 6, 3.

Symmetry with respect to an affine linear manifold: 6, 6.

Symmetric (matrix): 3, 1.

Symplectic (automorphism): 5, 3.

Symplectic (base): 5, 1.

Symplectic (group): 5, 3.

Symplectic (matrix): 5, 3.

Symplectic (transformation): 5, 3.

Tangent of an angle: 10, 3.

Tangent (linear variety) to a quadric: 6, exercises 23 and 25.

Pythagorean theorem: 6, 6.

Witt's theorem: 4, 3.

Totally isotropic (submodule): 4, 1.

Totally orthogonal (submodule) to a submodule: 1, 3.

Totally orthogonal (variety) to a linear variety: 6, 6.

Totally singular (submodule): 4, 1.

Orthogonal transformation: 6, 2.

Symplectic transformation: 5, 3.

Unitary transformation: 6, 2.

Trigonometric (function): 10, 3.

Type of a quadratic form: 8, 1.

Unitary (automorphism): 6, 2.

Unitary (group): 6, 2.

Unitary (special group): 6, 2.

Unitary (matrix): 6, 2.

Unitary (transformation): 6, 2.

Isotropic linear manifold: 6, 6.

Linear manifold orthogonal to a linear manifold: 6, 6.

Totally isotropic linear manifold: 6, 6.

Orthogonal linear manifolds: 6, 6.

Vector orthogonal to a linear manifold: 6, 6.

Witt (ring of): 8, 3.

Witt (decomposition of): 4, 2.

Witt (group of): 8, 2.

Witt (theorem of): 4, 3.

\chapter*{\texorpdfstring{Definitions of Chapter \Roman{chapter}}{Definitions of Chapter IX}}
\phantomsection
\addcontentsline{toc}{chapter}{\texorpdfstring{Definitions of Chapter \Roman{chapter}}{Definitions of Chapter IX}}

\section*{Sesquilinear forms:}

Let A be a ring, \(\xi \to \overline{\xi}\) an involutive antiautomorphism of A, that is, a bijection of A onto itself such that \(\overline{(\xi + \eta)} = \overline{\xi} + \overline{\eta}\) , \(\overline{(\xi\eta)} = \overline{\eta} . \overline{\xi}\) , \(\overline{\overline{\xi}} = \xi\) for all \(\xi, \eta\) in A. A sesquilinear form \(\Phi\) on a left A-module E is a map from E × E into A such that

\[
\begin{array}{c c} \Phi (x + x ^ {\prime}, y) = \Phi (x, y) + \Phi (x ^ {\prime}, y), & \Phi (x, y + y ^ {\prime}) = \Phi (x, y) + \Phi (x, y ^ {\prime}) \\ \Phi (\alpha x, y) = \alpha \Phi (x, y), & \Phi (x, \alpha y) = \Phi (x, y) \bar {\alpha} \end{array}
\]

for all \(\alpha\in\mathbf{A},x,x^{\prime},y,y^{\prime}\) in E.

Let \(\varepsilon\) be an element of the center of A. One says that a sesquilinear form \(\Phi\) on E is \(\varepsilon\)-hermitian if \(\Phi(y, x) = \varepsilon \overline{\Phi(x, y)}\) for all \(x \in E\) , \(y \in E\) ; if \(\varepsilon = 1\) (resp. \(\varepsilon = -1\) ) one says that \(\Phi\) is hermitian (resp. anti-hermitian).

When the antiautomorphism \(\xi\to\overline{\xi}\) is the identity (which implies that the ring A is commutative), the corresponding sesquilinear forms are the bilinear forms. One then says "symmetric" instead of "hermitian", and "antisymmetric" instead of "antihermitian" (cf.~Chap.~III). A bilinear form \(\Phi\) such that \(\Phi(x,x)=0\) is said to be alternating; it is then antisymmetric, and the converse is true when A is a field of characteristic \(\neq2\) (cf.~Chap.~III).

One says that a form \(\varepsilon\)-hermitian satisfies condition (T) if for all \(x \in E\) , there exists \(\lambda \in A\) such that \(\Phi(x, x) = \lambda + \varepsilon \overline{\lambda}\) . This condition is always satisfied when \(\varepsilon = 1\) and \(A\) is a field of characteristic \(\neq 2\) , or when \(\Phi\) is alternating.

\section*{Quadratic forms:}

Let A be a commutative ring. A quadratic form Q on an A-module E is a map from E into A such that \(\mathrm{Q}(\alpha x) = \alpha^{2}\mathrm{Q}(x)\) for \(\alpha \in A, x \in E\) , and that the map

\[
(x, y) \rightarrow \Phi (x, y) = Q (x + y) - Q (x) - Q (y)
\]

is a bilinear form (necessarily symmetric), said to be associated with the quadratic form Q. One has \(\Phi(x, x) = 2Q(x)\) ; conversely, for every bilinear form \(\Psi\) on E, \(x \to \Psi(x, x)\) is a quadratic form on E;

when A is a field of characteristic ≠ 2, quadratic forms and symmetric bilinear forms on E therefore correspond bijectively.

\section*{Orthogonal elements:}

Let \(\Phi\) be a form \(\varepsilon\)-hermitian on a left A-module E. One says that two elements \(x, y\) of E are orthogonal (for \(\Phi\) ) if \(\Phi(x, y) = 0\) ; this relation is symmetric in \(x\) and \(y\) . For every submodule M of E, the set of \(x \in E\) which are orthogonal to all elements of M is a submodule denoted \(\mathbf{M}^{\circ}\) and called the orthogonal of the submodule M. One says that \(\Phi\) is non-degenerate if \(\mathbf{E}^{\circ} = \{0\}\) . When E is a finite-dimensional vector space and \(\Phi\) is non-degenerate, one has codim \(\mathbf{M}^{\circ} = \dim \mathbf{M}\) and \(\mathbf{M}^{\circ\circ} = \mathbf{M}\) for every subspace M of E; moreover, for every pair of subspaces M, N of E, one has

\[
(\mathbf {M} + \mathbf {N}) ^ {0} = \mathbf {M} ^ {0} \cap \mathbf {N} ^ {0}, \quad (\mathbf {M} \cap \mathbf {N}) ^ {0} = \mathbf {M} ^ {0} + \mathbf {N} ^ {0}.
\]

Suppose the ring A is commutative, and let Q be a quadratic form on the A-module E, \(\Phi\) the bilinear form associated with Q. One says that two elements of E are orthogonal (for Q) if they are orthogonal for \(\Phi\) ; the orthogonal of a submodule M (for Q) is the orthogonal M \(^{o}\) of M for \(\Phi\) . One says that Q is non-degenerate if \(\Phi\) is non-degenerate.

\section*{Isotropic elements and singular elements:}

Let \(\Phi\) be a form \(\varepsilon\)-hermitian on a left A-module E. An element \(x \in E\) is said to be isotropic if \(\Phi(x, x) = 0\) ; a submodule M of E is said to be isotropic if \(\mathbf{M} \cap \mathbf{M}^{\circ} \neq \{0\}\) (in other words, if the restriction of \(\Phi\) to M is degenerate); M is said to be totally isotropic if \(\mathbf{M} \subset \mathbf{M}^{\circ}\) (in other words, if the restriction of \(\Phi\) to M is zero). For every submodule M of E, \(\mathbf{M} \cap \mathbf{M}^{\circ}\) is totally isotropic. If E is a finite-dimensional vector space and if \(\Phi\) is nondegenerate, the following three conditions are equivalent for a subspace M of E : \(1^{\circ}\) M is nonisotropic ; \(2^{\circ}\) M \(^{\circ}\) is nonisotropic ; \(3^{\circ}\) E is the direct sum of M and M \(^{\circ}\) . Suppose the ring A is commutative, and let Q be a quadratic form on an A-module E, \(\Phi\) the bilinear form associated with Q. An element \(x \in E\) is said to be singular if \(\mathrm{Q}(x) = 0\) ; a submodule M of E is said to be singular (resp. totally singular) if \(\mathbf{M} \cap \mathbf{M}^{\circ}\) contains a singular element \(\neq 0\) (resp. if the restriction of Q to M is zero). Every singular element (resp. singular submodule, totally singular submodule) is isotropic (resp. isotropic, totally isotropic); the converse is true when A is a field of characteristic \(\neq 2\) .
